\documentclass[11pt]{article}
\usepackage[margin=1in]{geometry}
\usepackage[T1]{fontenc}
\usepackage[utf8]{inputenc}
\usepackage{amsmath,amssymb,amsthm}
\usepackage[hidelinks]{hyperref}
\newtheorem{theorem}{Theorem}
\newtheorem{lemma}[theorem]{Lemma}
\newtheorem{corollary}[theorem]{Corollary}
\newcommand{\CPsh}{\operatorname{CPsh}}
\newcommand{\PSH}{\operatorname{PSH}_{\mathrm{cl}}}
\newcommand{\DD}{d'd''}
\title{Classical recognition, Shilov reconstruction, and upper control\\for generated analytic weights}
\author{Working paper; independent review pending}
\date{9 October 2026}
\begin{document}
\maketitle
\begin{abstract}
We apply the precise classical-psh foundations of Gubler--Rabinoff to the
minimal continuous CLD sheaf and its decreasing generated weights. This
gives tests by all analytic curves after all analytic field extensions,
the local maximum principle, and the affinoid Shilov maximum principle.
We spell out two consequences for an analytic torus quotient: local
weights are reconstructed from finite Shilov boundaries of arbitrarily
small affinoid neighborhoods, and convergence at all such boundary
points implies uniform upper control against a continuous candidate.
These are applications of retained source theorems. We do not claim a new
maximum principle or deduce compactness, envelopes, or the full
Monge--Ampere existence and uniqueness theorem from these consequences.
\end{abstract}

\section{Definitions and the exact source foundations}
Let $K$ be a complete non-Archimedean field and let $V$ be a good
$K$-analytic space. In the application $V$ is an ordinary open chart of
the analytic surface
\[
 K=\mathbb C((s)),\quad |s|=e^{-1/8},\qquad
 X=((\mathbb G_m)^2)^{\mathrm{an}}/
       \langle(s^8,2),(3,s^8)\rangle .
\]
The quotient and nonalgebraization are separate retained geometric
foundations. Here we use only its good open torus atlas. On a chart let
$x_i=-\log|z_i|$ and
\[
 \rho=\tfrac12(x_1^2+x_2^2),\qquad \theta=\DD\rho.
\]
The reference potential changes by an affine function of the unit
logarithms on an overlap.

The sheaf $\CPsh_V$ is the smallest sheaf of finite continuous real
functions containing ordinary CLD smooth functions with positive
$(1,1)$-Hessian and closed under uniform limits. The positivity is the
ordinary positive semidefinite cone in bidegree $(1,1)$.
Let $\mathcal C_\theta$ consist of finite continuous $u$ such that
$\rho+u\in\CPsh_V$ locally. A generated weight is a pointwise decreasing
limit of members of $\mathcal C_\theta$; the finite-energy generated class
adds the stated common lower energy bound. The results below do not need
the energy bound.

We use the exact arXiv:2506.13548v2 snapshot of Gubler--Rabinoff, dated
17 September 2025~\cite{GR}. Its definition of an analytic curve is a
separated analytic space of pure dimension one. Definition3.4 defines
subharmonicity on an arbitrary curve by all pullbacks from connected
smooth strictly analytic curves over nontrivially valued analytic field
extensions. Such a pullback must be either Thuillier-subharmonic or
identically $-\infty$. Thus the extended value is retained without
excluding identically infinite components.

Definition4.1 says that an upper semicontinuous function
$h:V\to[-\infty,\infty)$ is classically psh if, for every analytic
field extension $L/K$, every $L$-analytic curve $C$, and every actual
analytic morphism $f:C\to V_L$, its pullback
$h\circ\pi_{L/K}\circ f$ is subharmonic in Definition3.4. This definition
does not restrict to algebraic curves. Write $\PSH(V)$ for this class.

The following exact source results are retained:
\begin{enumerate}
\item The smooth-psh bridge in Section1.2: ordinary CLD smooth psh
functions are classically psh. The source reduces to smooth curves after
algebraically closed nontrivially valued extension and applies Wanner's
Theorem4.7. We retain this analytic-curve bridge as a named source
foundation. Its recursive local curve, model, and semistable background
is not newly proved here.
\item Proposition4.2: classical psh functions form a sheaf, are closed
under nonnegative linear combinations and finite maxima, and are closed
under infima of decreasing nets. They are compatible with all analytic
field extensions and all actual analytic morphisms. The converse
functorial statement in the source requires finite surjectivity; no such
converse is used below.
\item Theorem6.1: a classically psh function attaining a local maximum
at a point of the relative interior of an analytic space is locally
constant at that point.
\item Theorem6.4: a classically psh function on any $K$-affinoid space
attains its maximum at a Shilov point. The theorem includes possibly
trivially valued $K$ and general affinoid radii.
\end{enumerate}
The local maximum theorem uses the source's Theorem5.4, connectivity by
analytic curves and, for connected boundaryless spaces, by overconvergent
analytic curves. The global affinoid theorem uses its boundary
stratification, Lemma6.3. These are source foundations with their exact
analytic scope, rather than assertions inferred from a catalogue.

\section{Recognition of the whole continuous sheaf}
\begin{lemma}\label{uniform}
The finite continuous classical-psh functions form a sheaf closed under
uniform limits on an open set.
\end{lemma}
\begin{proof}
Sheafiness is Proposition4.2(1). Suppose $h_j$ converges uniformly to a
finite continuous $h$. Pass to a subsequence, still denoted $h_j$, with
$\|h_j-h\|_\infty\le2^{-j}$. Define
\[
 g_j=h_j+3\,2^{-j}.
\]
Adding a constant preserves classical psh. Moreover
\[
 g_{j+1}-g_j
 \le 2^{-(j+1)}+2^{-j}-3\,2^{-(j+1)}=0.
\]
Thus $g_j$ decreases pointwise to $h$. Proposition4.2(4) proves that
$h$ is classically psh. For locally uniform limits the same proof is
made on smaller opens and sheafiness glues the conclusion.
\end{proof}

\begin{theorem}\label{recognition}
Every member of $\CPsh_V$ is classically psh. Every decreasing generated
local weight $\rho+u$ is classically psh, with its extended $-\infty$
values retained. The assertion includes every continuous-sheaf stage,
all analytic curve tests, and all analytic field extensions.
\end{theorem}
\begin{proof}
Consider the sheaf of finite continuous members of $\PSH$.
The smooth bridge shows that it contains the smooth generators of
$\CPsh$. Lemma\ref{uniform} shows uniform closure.
The defining minimality gives $\CPsh\subset\PSH$ at all stages.
If $u_j\in\mathcal C_\theta$ decreases to $u$, the local functions
$\rho+u_j$ are classically psh and decrease to $\rho+u$.
Apply Proposition4.2(4) and then its field-extension and pullback rules.
\end{proof}
The minimality argument does not identify $\CPsh$ with its first
locally uniformly smooth-approximable stage. It proves one inclusion
into a different sheaf. The underlying classical-psh bridge is already
stated by Gubler--Rabinoff; no novelty is claimed for that bridge.

\begin{corollary}\label{maximum}
A generated local weight satisfies the local maximum principle at
relative interior points and the Shilov maximum principle on every
affinoid domain contained in the chart. In particular, for such a domain
$A$ with finite Shilov boundary $\Gamma_A$,
\[
 \max_{A}(\rho+u)=\max_{\Gamma_A}(\rho+u).
 \tag{1}
\]
If every boundary value is $-\infty$, the local weight is $-\infty$
throughout $A$.
\end{corollary}
\begin{proof}
Theorem\ref{recognition} and the exact source Theorems6.1 and6.4 apply.
Restriction to an affinoid domain is an actual analytic morphism, so
Proposition4.2(6) justifies that restriction. Upper semicontinuity gives
the maximum on compact $A$; the source theorem places it in $\Gamma_A$.
The case of an identically infinite function is allowed by its definition.
\end{proof}
It is $\rho+u$, rather than $u$ alone, that obeys (1). The reference
potential must remain in the boundary values.

\section{Reconstruction from arbitrarily small finite boundaries}
For a point $x\in V$, write $\mathcal A_x(V)$ for the affinoid
neighborhoods of $x$ contained in $V$, with $x$ in their ordinary
topological interior. The standard neighborhood property of good
analytic spaces supplies such a neighborhood inside every ordinary
open neighborhood of $x$. Let
\[
 D(V)=\bigcup_{A\subset V\text{ affinoid}}\Gamma_A.
\]
The boundaries in this union vary with $A$; this is not one fixed finite
set or one prescribed skeleton.

\begin{theorem}\label{reconstruction}
For every classically psh $q:V\to[-\infty,\infty)$,
\[
 q(x)=\inf_{A\in\mathcal A_x(V)}\ \max_{\xi\in\Gamma_A}q(\xi).
 \tag{2}
\]
Consequently, for every generated weight on a chart,
\[
 u(x)=\inf_{A\in\mathcal A_x(V)}
       \max_{\xi\in\Gamma_A}\bigl(\rho(\xi)+u(\xi)\bigr)-\rho(x).
 \tag{3}
\]
All these identities retain points where the weight is $-\infty$.
\end{theorem}
\begin{proof}
For each $A$ the Shilov maximum theorem gives
$q(x)\le\max_{\Gamma_A}q$. This is one inequality in (2).
If $q(x)$ is finite, for any $\epsilon>0$ upper semicontinuity gives an
ordinary open neighborhood $O$ of $x$ on which $q<q(x)+\epsilon$.
Choose $A\in\mathcal A_x(V)$ contained in $O$. Then
$\max_{\Gamma_A}q\le q(x)+\epsilon$.
Take the infimum and let $\epsilon$ decrease to zero.
If $q(x)=-\infty$, upper semicontinuity gives such an $O$ with
$q<B$ for every prescribed real $B$. The same argument shows that the
infimum is at most every $B$, hence is $-\infty$.
Apply (2) to $q=\rho+u$ to obtain (3).
\end{proof}

\begin{corollary}\label{order}
For classical-psh $q,r$ on $V$, if $q\le r$ on $D(V)$ then $q\le r$
on $V$. For two generated weights with the same reference potential,
their order, and hence equality, is determined by their values on
$D(V)$.
\end{corollary}
\begin{proof}
For each affinoid $A$ the assumed inequality on $\Gamma_A$ gives
$\max_{\Gamma_A}q\le\max_{\Gamma_A}r$.
Take infima in (2). For weights apply this to $\rho+u$ and $\rho+v$.
\end{proof}
Every nonempty ordinary open contains a nonempty affinoid domain, whose
Shilov boundary is nonempty. Thus $D(V)$ is dense. Nevertheless
Corollary\ref{order} is not merely the determination of a continuous
function by dense values: the weights may be discontinuous or have poles,
and (2) provides the needed upper-semicontinuous reconstruction.

\section{Upper control near a continuous candidate}
\begin{theorem}\label{Hartogs}
Let $(q_j)$ be a net of classical-psh functions on the good space $V$.
Let $h$ be a finite continuous real function on $V$, and assume
\[
 \limsup_j q_j(\xi)\le h(\xi)\qquad\text{for every }\xi\in D(V).
 \tag{4}
\]
Then, for every compact $Q\subset V$ and every $\epsilon>0$,
$q_j\le h+\epsilon$ on $Q$ for all sufficiently large $j$.
In particular
\[
 \limsup_j\sup_Q(q_j-h)\le0.
 \tag{5}
\]
\end{theorem}
\begin{proof}
Fix $\eta>0$. For each $x\in Q$ choose an affinoid neighborhood $A_x$
contained in $V$ on which the oscillation of $h$ is less than $\eta$.
This is possible by continuity and goodness. Finitely many ordinary
interiors of these neighborhoods cover $Q$; write them $A_1,\ldots,A_m$.
Their Shilov boundaries have finite union. By (4), a common late index
gives $q_j(\xi)\le h(\xi)+\eta$ at every point of that finite union.
For $y\in Q\cap A_i$, the affinoid maximum principle now gives
\[
 q_j(y)\le\max_{\Gamma_{A_i}}q_j
 \le\max_{\Gamma_{A_i}}h+\eta
 \le h(y)+2\eta.
\]
The common late index exists for a net because its index set is directed.
Take $2\eta<\epsilon$ to obtain the assertion and then (5).
\end{proof}

\begin{corollary}\label{weighted}
Let $(u_j)$ be generated $\theta$-psh weights on the compact quotient
$X$, and let $f$ be finite continuous. Suppose that on every reference
chart the boundary-value condition
$\limsup_j u_j(\xi)\le f(\xi)$ holds at all its affinoid Shilov points.
Then $u_j\le f+\epsilon$ on $X$ for every $\epsilon>0$ and every
sufficiently late $j$.
\end{corollary}
\begin{proof}
On each chart apply Theorem\ref{Hartogs} to $q_j=\rho+u_j$ and
$h=\rho+f$. Compactness and the good open atlas give a finite collection
of compact affinoid neighborhoods contained in the charts and covering
$X$. Apply the theorem on these compact sets and take a common late
index. Subtract the corresponding reference potential on each chart.
\end{proof}
This conclusion controls only the upper errors. It neither asserts
uniform convergence nor bounds the negative tails of an energy sublevel.

\section{The remaining compactness and equation gates}
Equations (2)--(5) replace a previously missing maximum and determination
input for the full generated class. The ordinary analytic Shilov points
here are not an invariant-measure support restriction. A Monge--Ampere
measure is not assigned by these boundary evaluations.

These facts still do not supply a compactness theorem for normalized
finite-energy potentials. A limit of their values on $D(V)$ need not
already be the boundary values of a generated weight. Nor is an arbitrary
pointwise limit automatically upper semicontinuous. The exact source
Theorem6.5 assumes that a pointwise net limit is upper semicontinuous
before concluding it is classical psh; it cannot remove this requirement.
Reconstruction from genuine psh values in (2) is not a construction of
every prescribed boundary-value limit.

No converse identification with all intrinsic psh weights, continuous
obstacle envelope theorem, orthogonality, uniform integrability of all
finite-dual measures on energy sublevels, arbitrary-support equation
existence, or full uniqueness follows here. Original P16.7 remains open.
The proposed applications and this working paper await independent review.

\section{Source version and audit boundary}
The complete applied definitions and immediate proofs in the retained
Gubler--Rabinoff v2 text were read for Proposition4.2 and
Theorems5.4,6.1,6.4,6.5, together with the actual boundary stratification
Lemma6.3 and the relevant curve maximum, base-change, and sheaf proofs.
The exact smooth bridge is stated and argued in Section1.2.
The recursive Berkovich, Ducros, Temkin, Thuillier, and Wanner proof trees
remain named source foundations, not entire freshly audited proof trees.
In particular we use no unsupported arbitrary-current pullback.

The actual original PDF pages3,33,34 were also inspected. Page3 contains
the analytic-curve bridge, and pages33--34 contain the full statement and
immediate proof of Theorem6.4 and the pointwise-limit theorem.
The extracted and printed proof of Theorem6.4 says ``since $h$ is continuous'' when
claiming that the maximum set is closed. The theorem's stated hypothesis
is upper semicontinuity; at the attained supremum that condition already
makes the maximum set closed. Our proofs use exactly upper
semicontinuity. The literal source text is retained without silently
altering its words. The snapshot is not identified byte for byte with
any later publication, and no complete new reading of the entire paper
or of its recursive dependencies is claimed.

\begin{thebibliography}{9}
\bibitem{GR} W. Gubler and J. Rabinoff,
\emph{Classically psh and pluriharmonic functions on Berkovich spaces},
arXiv:2506.13548v2,17 September2025,
\url{https://arxiv.org/abs/2506.13548v2}.
\bibitem{Wanner} V. Wanner,
\emph{Comparison of two notions of subharmonicity on non-Archimedean
curves}, arXiv:1801.04713v3,8 November2018.
Published in Math. Z.293 (2019),443--474.
\bibitem{CLD} A. Chambert-Loir and A. Ducros,
\emph{Formes differentielles reelles et courants sur les espaces de
Berkovich}, arXiv:1204.6277v2,25 July2025.
The minimal continuous-sheaf convention is Definition8.1.7.
\end{thebibliography}
\end{document}
